\section{L'impasse ergonomique}
Si l'infrascture technique de CujasNum pose d'importants défis en termes de sécurité, la limite se révèle dans son fonctionnement quotidien. La chaîne de numérisation actuelle de la bibliothèque numérique illustre parfaitement l'approche décrite par Elina Leblanc dans sa thèse la qualifiant d'une \enquote{approche purement descendante(top-down), axée sur les contraintes du système plutôt que sur l'activité des agents}.

Elina Leblanc s'appuye sur les critères d'Aaron Schmidt, l'évaluation d'un service de bibliothèque numérique présente \enquote{un guide pour la CCU des bibliothèques analogiques [...]sous un angle utilisateur}. Il mentionne ainsi l'utilité d'un service d'une bibliothèque numérique en disant que celui-ci doit être \enquote{utilisable et désirable}. Aaron Schmidt l'explique ainsi : \enquote{Un service utile doit répondre à un besoin et avoir une importance aux yeux des utilisateurs pour qu'ils ne l'utilisent pas de manière ponctuelle, mais sur la longue durée. Un service utilisable garantit aux utilisateurs un accès aisé et rapide aux données. Enfin, la désirabilité crée un lien entre l'utilisateur et un service, en donnant envie au premier d'utiliser le deuxième et en suscitant un besoin}. Cela garantit à l'utilisateur une expérience optimale lorsqu'ils utilisent la bibliothèque numérique d'une institution patrimoniale. Hors, à la bibliothèque Cujas, l'interface n'est pas conçue pour s'adapter à l'humain et à la manière dont celui-ci est utilisé les sites web aujourd'hui. L'utilisateur doit s'adapter lui aux contraintes rigides de l'interface et de la manière dont les services lui sont présentés. 

Ce manque d'ergonomie se traduit par un décalage dans la production des ouvrages numériques. Si le contrôle qualité d'un ouvrage requiert environ \enquote{quinze minutes en moyenne si tout va bien}, la moindre anomalie structurelle ou typographique fait basculer le traitement dans une autre dimension nécessitant à minima \enquote{trois jours} de manipulations afin de corriger l'erreur détectée et ensuite pouvoir remettre en ligne le document. Cette manière de travailler est devenue un gouffre chronophage et la rigidité qu'implique ce traitement est devenu un carcan dont la bibliothèque Cujas ne parvient pas à se défaire.

\subsection{La pénibilité des traitements manuels et de l'encodage XML}
Sous CujasNum, la production n'est pas automatisé à toutes les échelles, mais plutôt un \enquote{bricolage permanent} afin de permettre une autonomie des équipes tout en conservant la supervision par un regard humain des tâches effectuées. Ces caractéristiques de travail illustre les limites de \enquote{la pensée documentaire du deuxième âge du numérique}. 

Le premier goulot d'étrangelement de cette chaîne réside dans l'obligation de manipuler plusieurs formats et de les maintenir simultanément pour une même donnée structurelle. En effet, c'est le cas pour leur table des matières (TDM) d'un ouvrage, qui sont retranscrite manuellement afin de reconstituer intellectuellement l'ouvrage. Cette TDM est d'ailleurs essentiel à l'utilisateur pour naviguer à l'intérieur d'un document numérisé. Cette TDM est saisie dans un fichier Excel. Par la suite, ce même fichier Excel est utilisé pour générer celle-ci sous le format CSV.

A cela s'ajoute l'obligation de respecter une présentation précise : le premier onglet du tableau Excel s'intitule obligatoirement \enquote{Feuille 1}, qui est la nomination standardisée. Les agents doivent alors vérifier avec le livre en main, afin de s'assurer de la concordance entre les pages numérisées et les pages réelles, ainsi qu'avec ce qui a été indiqué dans le bordereau de commande. Les agents doivent \enquote{vérifier qu'il y a bien une image pour chaque page} ainsi que le nommage du fichier image nommé avec \enquote{le code-barres suivi du numéro d'image et de l'extension}. 

Le coût cognitif le plus important des équipes de Cujas apparaît au moment de la validation sémantique et syntaxique du code XML. Les agents effectuent ce contrôle manuellement afin de s'assurer de la conformité de l'encodage avant de pouvoir l'intégrer dans la bibliothèque numérique et de l'envoyer au CINES sans que celui-ci ne soit rejeter. Pour éviter ce rejet, cela passe par plusieurs vérifications notamment avec l'utilisation des guillemets. Le technicien doit s'assurer que ce soit les guillemets anglais qui sont utilisés et non français, car dans le cas où ce sont les guillemets français, ils créent des problèmes d'interprétation par le parseur XML. Ils sont interprétés comme des perturbations de structuration XML qui sont déjà utilisé pour les balises de celui-ci entraînant donc un conflit. Pour le résoudre, il a été décidé d'utiliser les guillemets anglais, mais n'étant pas à l'abri d'un oubli, c'est une étape fastidieuse mais indispensable à réaliser afin d'éviter de futur problèmes au moment d'exécuter la procédure de mise en publication. Parmi les caractères spéciaux, d'autres sont également interdit et requièrent une intervention humaine ce sont les esperluettes \&. Elles doivent être répérés et encodés manuellement sous la forme d'une entité XML en utilisant par exemple \&amp;.

En plus de ces corrections, le technicien en charge du contrôle doit également s'assurer de la conformité du schéma XML car le CINES n'accepte plus les liens commençant par \enquote{http} mais privilégient ceux qui utilisent \enquote{https}. Une attention doit être également portée sur les sauts descendants incohérents. \enquote{Un agent doit s'assurer qu'un document ne passe pas directement d'un niveau hiérarchique 1 (Chapitre) à un niveau 3 (sous-section) sans déclarer formellement le niveau 2 intermédiaire}. Ces règles de logique hiérarchique s'applique à la structuration des tables des matières. 

Toutes cces micro-corrections, indispensables pour empêcher le blocage lors de l'affichage ou du versement des fichiers reposent sur des interventions humaines précises et mettent en avant l'idée d'un \enquote{bricolage manuel permanent} afin de s'assurer de la conformité des notices. L'infrastructure technique ne propose pour le moment aucune interface d'aide à la saisie ou de correction automatique afin d'éviter aux équipes ce travail méticuleux et chronophage. 

Ces micro-intervention créent une défaillance ergonomique entraînant une asymétrie temporelle dans les différentes étapes de la chaîne de numérisation. En temps normal, le contrôle qualité d'un ouvrage dure environ \enquote{quinze minutes en moyenne si tout va bien}. En revanche, lorsque les scripts détectent des erreurs cela peut prendre proportion inadapté contenue de l'erreur initiale. Habituellement, la correction d'une erreur simple devrait être proportionnelle à sa gravité. Or, à Cujas, cela n'est pas le cas, un déséquilibre important c'est crée entre la légèreté d'une correction, par exemple l'oubli de remplacer les guillemets français par ceux anglais, et la lourdeur du processus de correction. Actuellement l'interface n'offre aucune possibilité de faire une modification rapide \textit{in situ} en évitant d'avoir à relancer toute la procédure pour une coquille. 

En effet, l'apparition d'une métadonnée erronnée ou d'un bug oblige en l'état, l'agent à \enquote{extraire le document, à corriger manuellement les fichiers METS ou Dublin Core sous Oxygen, puis à ré-uploader le lot via un protocole SFTP sécurisé (WinSCP) dont le traitement par lot n'est planifié que durant la nuit}. La rectification d'une coquille ou d'un lien brisé peut nécessiter jusqu'à trois jours de traitement minimum avant d'être visible en ligne pour l'usager.  

\subsection{La rigidité d'XTF face aux formats non monographiques}
À l'origine l'architecture d'XTF a été pensée exclusivement pour un seul type de document -- la monographie -- et qu'elle se révèle incapabe de s'adapter à d'autres types d'arborescences, comme les périodiques. La particularité d'XTF repose sur une arborescence dite \enquote{à plat} ainsi il n'est pas impossible d'ajouter d'autres types de documents dans la bibliothèque numérique excepté que l'arborescence des documents ne sera alors pas respecté. Cela représente un problème majeur pour les équipes de Cujas qui se trouvent cantonner aux monographies sans avoir la possibilité d'ajouter leurs autres collections. 

Une des nombreuses contraintes techniques d'XTF réside dans l'obligation d'avoir le format PDF afin de mettre le document en ligne. Le problème majeur vient surtout que le système a été conçu comme \enquote{un simple répertoire PDF} calqué sur la structure physique d'un livre (volume et pages). Cela rend donc difficile l'ajout de tout autre types de documents ne respectant pas ce principe. Pour rendre la consultation des images plus agréables, une visionneuse JPEG a été implémenté en 2013 pour permettre à l'usager de feuilleter le document de façon dynamique. 

Contrairement à une monographie, un périodique exige une architecture d'accès d'au moins de trois niveau, c'est-à-dire le titre de la revue, suivie de l'année de publication et du numéro du volume ou du fascicule, pour que la consultation soit pertinente. XTF est incapable de reproduire cette structure hiérarchique c'est donc la raison principale qui a poussé la bibliothèque Cujas à mettre de côté ses collections ne pouvant les intégrer sans \enquote{dénaturer} le type de document. Aujourd'hui, on observe sur certaines bibliothèques numériques possédant des périodiques que la granularité peut aller jusqu'à l'article même. XTF s'arrête initialement au volume et ne permet pas une granularité aussi précise, allant donc à l'encontre des pratiques de recherche actuelle. Faute de pouvoir diffuser ces revues dans sa bibliothèque numérique, Cujas a fait le choix de faire numériser ces volumes via un partenariat avec Persée. Pour signaler ces numérisations, la bibliothèque se contente uniquement de créer des notices de redirection indiquant aux chercheurs l'existance de ces documents où est-ce qu'il peut les trouver, afin de pouvoir les consulter. 

La bibliothèqe Cujas rencontre le même problème avec son fonds iconographiques qui ne possédant pas de structure intellectuelle répliquable sous la forme d'une table des matières. Si les équipes de Cujas souhaitent mettre en ligne ce fonds. Cela peut être effectuée \enquote{de manière détournée en créant de PDF}. Ils doivent alors créer \enquote{artificiellement une table des matières et une table de concordance}, afin de pouvoir mettre en ligne ces documents. Cependant, c'est au niveau de l'archivage pérenne que le blocage risque de se faire suite à ce type de bricolage. En effet, cela obligerait la bibliothèque Cujas à \enquote{passer par une redéfinition complexe du plan d'archivage avec le CINES}. 

Un autre problème réside dans l'indexation utilisée par la bibliothèque numérique actuelle, l'indexation Hauville. En effet, celle-ci ne permet pas le filtrage via différents types de documents. L'indexation n'a a été structuré qu'autour de la monographie ne permettant pas de croiser les informations. L'impossible d'ajouter d'autres types de documents met également en luminière l'invisibilisation auquel ils se trouveraient confrontés lors d'une recherche plein-texte. Ces documents ne pourraient ressortir de manière évidente ne pouvant posséder d'OCR. Les résultats apparaisseront donc de façon très restreinte et dans les dernières pages problablement. L'obsolescence du logiciel XTF entraîne donc \textit{de facto} un filtrage empêchant la bibliothèque Cujas de pouvoir mettre en avant la richesse et la diversité de ses collections. Elle ne peut pas non plus ajouter de valorisation des documents afin de justifier les choix effectués ou de mettre en avant une collection ou un document dans un article ou dans une exposition virtuelle. 
\medskip

Le bilan de CujasNum met en avant les points forts du projet ainsi que les points faibles au moment de la réalisation de celui-ci mais également ceux ayant émergé par la suite avec l'évolution des technoloqies numériques ainsi que des pratiques du Web. D'abord, pionnière dans le monde des bibliothèques numériques, la bibliothèque Cujas s'est trouvé progressivement isolée ne parvenant pas suivre l'évolution rapide liée à l'émergence de l'internet entraînant les pratiques documentaires et les pratiques d'Internet dans un nouvel âge. Fâce à ce constat, l'ajout d'une surcouche WordPress ne s'aurait être la solution comme cela avait été envisagé dans un premier temps. Pour répondre aux attentes contemporaines des chercheurs et pour s'adapter aux mutations nationales des systèmes de gestion des bibliothèques (SGB) impose une rupture dans les pratiques actuelles. Il leur est donc nécessaire de penser une interface plus souple, acceptant différents types de documents se basant sur les pratiques actuelles de la recherche afin de répondre aux nouveaux enjeux des utilisateurs. La bibliothèque numérique doit permettre à sa bibliothèque d'utiliser les standards du Web de données afin de rendre sa plateforme davantage interopérable et prête à s'intégrer dans de grands réseaux scientifiques nationaux.